\section{The results, read table by table}\label{ch:results}

This chapter goes through the paper's results in the order of its results section. For each table or figure it says
which question it answers, how to read it, what it shows and what it does not. Unless a caption says otherwise, a cell
is the mean $\pm$ standard deviation over seeds 3--5, a difference is FedECS-LT minus the other method, paired over
the same seeds, and ``significant'' means $p<0.05$ after Holm's correction (\cref{sec:s-stats-design}). The noise
floor is \R{repeat/bal/floor} points without the readout and \R{repeat/er/floor} with it.

Two labels recur. \emph{FedECS-LT} is the full method, ending with the exact readout. \emph{FedECS-LT (heads only)} is
the same run before the readout, the heads as Stage III left them. Every other method appears \emph{without the
readout} in the main tables, as its authors defined it; in \cref{sec:s-same-readout} the readout is added to the
methods that can take it.

\subsection{The main comparison}\label{sec:s-main}

\paragraph{The question} Does FedECS-LT raise balanced accuracy over every published method when clients have
different models and the training set has a long tail? This is RQ2.

\paragraph{How to read the tables} \Cref{tab:s-main} is the paper's main table: one row per method, one column per
benchmark, the best mean in each column in bold. \Cref{tab:s-paired} is its paired-difference table: each cell is FedECS-LT
minus the row's method, with the 95\% interval in brackets, bold where the difference is significant. A difference
whose interval excludes zero but is not bold passed the uncorrected test and failed Holm's correction.

\begin{table}[t]
\centering\small
\caption{Balanced accuracy (\%), the paper's main table. Look at the first three rows: the full method, its heads before
the readout, and GeFL-F; then at the strongest other method in each column.}
\label{tab:s-main}
\slottab{main_bal}{}
\end{table}

\begin{table}[t]
\centering\small
\caption{FedECS-LT minus each method (points), paired over seeds 3--5, with the 95\% $t$ interval; bold: significant
after Holm's correction. The paper's paired-difference table. Look at the smallest difference in each column:
it belongs to the strongest other method.}
\label{tab:s-paired}
\slottab{main_paired}{}
\end{table}

\paragraph{The numbers} On CIFAR-10-LT FedECS-LT reaches \R{main/cifar10/FedECS-LT/bal}, against
\R{main/cifar10/strongest/bal} for the strongest other method, \R{main/cifar10/strongest}: a margin of
\R{main/cifar10/margin} points with Holm-corrected $p=\R{main/cifar10/strongest_p}$. On CIFAR-100-LT it reaches
\R{main/cifar100/FedECS-LT/bal} against \R{main/cifar100/strongest/bal} for \R{main/cifar100/strongest}, a margin of
\R{main/cifar100/margin} with $p=\R{main/cifar100/strongest_p}$. Against GeFL-F the differences are
\R{main/cifar10/diff/GeFL-F} \R{main/cifar10/ci/GeFL-F} and \R{main/cifar100/diff/GeFL-F}
\R{main/cifar100/ci/GeFL-F}. On CIFAR-100-LT no class has zero recall for all ten models of FedECS-LT
(\R{main/cifar100/FedECS-LT/nzero}), while \R{main/cifar100/strongest_nzero} classes on average are recognised by none
of the models of \R{main/cifar100/strongest}.

\paragraph{The heads alone} Before the readout the heads of FedECS-LT reach
\R{main/cifar10/FedECS-LT (heads only)/bal} and \R{main/cifar100/FedECS-LT (heads only)/bal}. That is
\R{main/cifar10/heads_margin} points ahead of \R{main/cifar10/heads_strongest} on CIFAR-10-LT, with
$p=\R{main/cifar10/heads_strongest_p}$, not significant, and \R{main/cifar100/heads_margin} ahead on CIFAR-100-LT, with
$p=\R{main/cifar100/heads_strongest_p}$, significant. The readout itself adds \R{main/cifar10/FedECS-LT/gain_er} and
\R{main/cifar100/FedECS-LT/gain_er} points to the heads. Of the \R{main/cifar10/margin}-point margin on CIFAR-10-LT,
therefore, about a third comes from the heads and two thirds from the readout.

\paragraph{What it shows} The measured statement is narrow: the full method, with its readout, is above every
published method of this setting as that method was published, by a large and significant margin on both benchmarks.
The heads alone are also ahead, by less, and the CIFAR-10-LT part of that is not resolved by three seeds.

\paragraph{What it does not show} It does not show that the generator side of FedECS-LT is what helps. Every other
method in this table lacks the readout, and the readout is a known kind of classifier (\cref{sec:s-readout}) that can be
added to most of them. The comparison that answers this follows in \cref{sec:s-same-readout}.

\subsection{The other model-heterogeneous methods}

\paragraph{The question} Do the methods that share prototypes, a generator of embeddings or a shared sub-network do
better than the generator-based ones under a long tail?

\paragraph{The answer} No. In \cref{tab:s-main} FedProto, FedTGP, FedGen, FedGH, LG-FedAvg and FedMRL are close to or
below Local. None of them gives an architecture any information about the classes its own clients lack: a prototype of
class 9 is only useful to a head that already predicts class 9. FedProto, FedTGP and FedGen also decline during
training. They are best in the first rounds and lose \R{main/cifar10/FedProto/gefl_drop},
\R{main/cifar10/FedTGP/gefl_drop} and \R{main/cifar10/FedGen/gefl_drop} points by the last round on CIFAR-10-LT. The
heads of FedECS-LT decline too, by \R{main/cifar10/FedECS-LT (heads only)/gefl_drop} points on CIFAR-10-LT. The paper's
appendix gives every method's best round, the measure GeFL reports; the readout is computed once, after the last round,
so it has no best round.

\paragraph{What it does not show} These methods run as ports onto the frozen extractor with settings from their
library, meant for end-to-end training (\cref{ch:design}, threats). Their low numbers may partly be the port. The paper
did not run a check of the ports.

\subsection{Accuracy on the clients' own class mixes}

\paragraph{The question} Balanced accuracy weights every class equally. What happens to the accuracy each client
actually sees, on test images drawn with its own class frequencies?

\paragraph{The numbers} FedECS-LT reaches \R{main/cifar10/FedECS-LT/local} on CIFAR-10-LT and
\R{main/cifar100/FedECS-LT/local} on CIFAR-100-LT, and its heads alone \R{main/cifar10/FedECS-LT (heads only)/local}
and \R{main/cifar100/FedECS-LT (heads only)/local}. GeFL-F reaches \R{main/cifar10/GeFL-F/local} on CIFAR-10-LT and
Local \R{main/cifar10/Local/local}. FedECS-LT is below \R{main/cifar10/local_n_above} of the
\R{main/cifar10/local_n_other} other federated methods on CIFAR-10-LT and \R{main/cifar100/local_n_above} of
\R{main/cifar100/local_n_other} on CIFAR-100-LT.

\paragraph{Why} This is the price of balanced accuracy and of the readout, and it follows from their definitions. A
client whose test images are mostly class 0 is best served by a last layer that favours class 0. Logit adjustment and
the readout's equal class priors remove exactly that preference. The readout, which treats all classes as equally
likely, lowers this metric by about nine points on CIFAR-10-LT. A deployment that cares about each client's own
accuracy could add each client's own log-prior back to the readout's biases at prediction time, the reverse of
\cref{sec:s-lastlayer}; the paper did not test this.

\subsection{The same readout for every method}\label{sec:s-same-readout}

\paragraph{The question} How much of FedECS-LT's margin remains when every method of the GeFL-F pipeline gets the same
readout? This is RQ3, and its answer decides how the paper's claim must be worded.

\paragraph{Why this is the fair comparison of the bodies} With the readout, every last layer is replaced by the same
estimator, computed from the same kind of pooled statistics, through the same procedure. What still differs between two
methods is the body that produced the embeddings. So the difference with the readout measures what FedECS-LT's
other components (the anchored generator, the moment correction, consensus soft labels, logit adjustment and ten
synthetic epochs) do to the bodies, read through one linear classifier. The comparison without the readout measures
the methods as published. A reader needs both: the first says whether FedECS-LT beats what exists; the second says
whether its parts beyond a known classifier contribute.

\begin{table}[t]
\centering\small
\caption{Balanced accuracy (\%) with every last layer set by the exact readout, the paper's readout table. Methods of the
GeFL-F pipeline only. Compare each row with its value in \cref{tab:s-main}: the readout raises every method.}
\label{tab:s-er}
\slottab{main_er}{}
\end{table}

\begin{table}[t]
\centering\small
\caption{FedECS-LT minus each method, both with the readout (points), paired over seeds 3--5, with the 95\% $t$
interval; bold: significant after Holm's correction. From the paper's appendix. Look at whether the intervals contain
zero.}
\label{tab:s-er-paired}
\slottab{main_er_paired}{}
\end{table}

\paragraph{The numbers} With the readout (\cref{tab:s-er,tab:s-er-paired}) FedECS-LT is \R{main/cifar10/er_diff/GeFL-F}
points \R{main/cifar10/er_ci/GeFL-F} ahead of GeFL-F on CIFAR-10-LT and \R{main/cifar100/er_diff/GeFL-F}
\R{main/cifar100/er_ci/GeFL-F} on CIFAR-100-LT. Against the strongest other method with the readout,
\R{main/cifar10/er_strongest} on CIFAR-10-LT and \R{main/cifar100/er_strongest} on CIFAR-100-LT, the margins are
\R{main/cifar10/er_margin} \R{main/cifar10/er_strongest_ci} and \R{main/cifar100/er_margin}
\R{main/cifar100/er_strongest_ci}, with Holm-corrected $p=\R{main/cifar10/er_strongest_p}$ and
$\R{main/cifar100/er_strongest_p}$. On CIFAR-100-LT the uncorrected interval excludes zero but the corrected test does
not reach $0.05$. The differences to GeFL-F and FedDF on CIFAR-100-LT stay significant
($p=\R{main/cifar100/er_p/GeFL-F}$ and $\R{main/cifar100/er_p/FedDF}$).

\paragraph{What it shows} Three things. First, the readout raises every method by a large amount: Balanced Softmax
gains \R{main/cifar10/Balanced Softmax/gain_er} points on CIFAR-10-LT, GeFL-F \R{main/cifar10/GeFL-F/gain_er}. Second,
once every method has it, the methods are close. Third, Balanced Softmax with the readout reaches
\R{main/cifar10/Balanced Softmax/er_bal} on CIFAR-10-LT. With one client per architecture, Balanced Softmax trains every
head on its own client's data with no generator and no exchange after Stage I. So a body trained alone, plus the
federated readout, comes within \R{main/cifar10/er_diff/Balanced Softmax} points of the full method. Most of the margin
of \cref{sec:s-main} is the readout's.

\paragraph{What it does not show} It does not show that the other components do nothing: a gain of 1.5 points with
three seeds and a noise floor of \R{repeat/er/floor} is neither shown nor ruled out. It also does not show the
components' effect without the readout, which is the subject of the pending ablations (\cref{sec:s-pending}). And
because the readout is one fixed linear classifier, a body that is better for another last layer would not show
it here.

\subsection{Rare classes}

\paragraph{The question} Where among the classes does the gain sit? If FedECS-LT addresses the long tail, the gain
should be largest on the rare classes.

\paragraph{How to read it} \Cref{tab:s-tail} gives the mean recall of the tail classes (the 40\% least frequent:
classes 6--9 on CIFAR-10-LT, 60--99 on CIFAR-100-LT). \Cref{fig:s-recall-size} plots each class's recall against its
number of training images on CIFAR-100-LT.

\begin{table}[t]
\centering\small
\caption{Recall (\%) of the tail classes, the 40\% least frequent; FedECS-LT with its readout, other methods without
it. From the paper's tail table. Compare FedECS-LT with its heads only and with GeFL-F.}
\label{tab:s-tail}
\slottab{main_tail}{}
\end{table}

\begin{figure}[t]
\centering
\slotfig{fig_recall_size}{col}{Recall of every class against its training size, CIFAR-100-LT.}
\caption{Recall of each class of CIFAR-100-LT against its number of training images in the federation (logarithmic
axis); each point is one class, mean over seeds 3--5. Look at the left side of the plot, the rarest classes, where the
methods separate.}
\label{fig:s-recall-size}
\end{figure}

\paragraph{The numbers} Against GeFL-F on CIFAR-10-LT, FedECS-LT gains \R{main/cifar10/tail_diff/GeFL-F} points on
the tail classes, \R{main/cifar10/med_diff/GeFL-F} on the medium ones and \R{main/cifar10/head_diff/GeFL-F} on the head
classes. Its tail recall is \R{main/cifar10/FedECS-LT/tail} against \R{main/cifar10/GeFL-F/tail}. Against FedDF it
gains \R{main/cifar10/tail_diff/FedDF} on the tail and loses \R{main/cifar10/head_diff/FedDF} on the head classes. On
CIFAR-100-LT the gain over GeFL-F is spread over the three groups (head \R{main/cifar100/head_diff/GeFL-F}, medium
\R{main/cifar100/med_diff/GeFL-F}, tail \R{main/cifar100/tail_diff/GeFL-F}). There the tail recall is
\R{main/cifar100/FedECS-LT/tail} against \R{main/cifar100/GeFL-F/tail}, and the worst-class recall
\R{main/cifar100/FedECS-LT/worst} against \R{main/cifar100/GeFL-F/worst}: some class is still lost on CIFAR-100-LT.

\paragraph{With the readout for every method} The tail gain over the loss-only methods disappears. FedECS-LT minus
Balanced Softmax, both with the readout, is \R{main/cifar10/ertail_diff/Balanced Softmax} points of tail recall on
CIFAR-10-LT and \R{main/cifar100/ertail_diff/Balanced Softmax} on CIFAR-100-LT. Against GeFL-F with the readout a tail
gain remains, \R{main/cifar10/ertail_diff/GeFL-F} \R{main/cifar10/ertail_ci/GeFL-F} on CIFAR-10-LT and
\R{main/cifar100/ertail_diff/GeFL-F} \R{main/cifar100/ertail_ci/GeFL-F} on CIFAR-100-LT, significant on CIFAR-100-LT
($p=\R{main/cifar100/ertail_p/GeFL-F}$).

\paragraph{What it shows and does not} The long-tail gain is real and sits where it should, on the rare classes; its
trade is head-class recall against FedDF. The same-readout numbers say where it comes from: from the last layer. The
part that remains against GeFL-F is consistent with the generator side helping the bodies of GeFL-F's own pipeline,
but against a body trained without any generator it is zero.

\subsection{Natural class imbalance}\label{sec:s-natural}

\paragraph{The question} Does the result carry over to datasets whose imbalance is their own, which played no part in
the development of the method? This is the first half of RQ4, and it is the cleanest test in the paper
(\cref{sec:s-development}).

\begin{table}[t]
\centering\small
\caption{Balanced accuracy (\%) on the six naturally imbalanced datasets, the paper's natural-imbalance table; FedECS-LT with its
readout, other methods without it. A value without $\pm$ is from one seed. Look at the gap between the first row and
the rest.}
\label{tab:s-natural}
\slottab{nat_bal}{}
\end{table}

\begin{table}[t]
\centering\small
\caption{The same datasets with the exact readout for every method of the GeFL-F pipeline, from the paper's appendix.
Look at how close the rows are.}
\label{tab:s-natural-er}
\slottab{nat_er}{}
\end{table}

\paragraph{The numbers} Without the readout for the others (\cref{tab:s-natural}), FedECS-LT has the highest mean on
\R{nat/nbest_of} datasets. Its margin over the strongest other method runs from \R{nat/margin_lo}
(\R{nat/margin_lo_ds}) to \R{nat/margin_hi} points (\R{nat/margin_hi_ds}), significant on \R{nat/n_sig} of the
\R{nat/n_sig_of} datasets with three seeds; \R{nat/one_seed} has one seed, so no test. With the readout for every method
of the pipeline (\cref{tab:s-natural-er}), the margins run from \R{nat/er_margin_lo} (\R{nat/er_margin_lo_ds}, against
\R{nat/er_margin_lo_who}) to \R{nat/er_margin_hi} (\R{nat/er_margin_hi_ds}). None is significant after Holm's
correction (the smallest $p$, on OrganAMNIST, is \R{nat/organamnist/er_strongest_p}). All but the BloodMNIST margin
exceed the noise floor of \R{repeat/er/floor} points, which is necessary but not sufficient for a resolved difference.
On BloodMNIST, FedDF with the readout has the higher mean. The readout changes FedECS-LT itself by
\R{nat/pathmnist/FedECS-LT/gain_er} points on PathMNIST and by \R{nat/organamnist/FedECS-LT/gain_er} on OrganAMNIST,
where it slightly lowers the heads' accuracy. On \R{gate/rows} the gate chose the generator with class rows.

\paragraph{What it shows} On data the development never saw, the pattern of the CIFAR benchmarks repeats: a large
margin as published, a small and unresolved one once every method has the readout. That the pattern repeats is evidence
that the CIFAR result is not an artefact of tuning. That the same-readout margin stays small is evidence that the
readout carries the result here too.

\paragraph{What it does not show} On OrganAMNIST, BloodMNIST and PathMNIST the realised imbalance ratio is 5 or less
(\cref{ch:design}), so these three test label skew more than a long tail.

\subsection{One shared model}\label{sec:s-shared}

\paragraph{The question} When all clients can run the same model, how does FedECS-LT compare with the federated
long-tail methods? This is the second half of RQ4. It is a test of scope, outside the setting the method was built for.

\begin{table}[t]
\centering\small
\caption{Balanced accuracy (\%) with one ResNet-8 shared by all clients (FedLTLib protocol), imbalance ratio 100; the
paper's shared-model table. A value without $\pm$ is from one seed. Look at how many rows are above the first one.}
\label{tab:s-shared}
\slottab{t2_bal}{}
\end{table}

\paragraph{The numbers} \R{t2/cifar10-r100/above_list} are above FedECS-LT on CIFAR-10-LT, and
\R{t2/cifar100-r100/above_list} on CIFAR-100-LT (\cref{tab:s-shared}). FedECS-LT reaches \R{t2/cifar10-r100/FedECS-LT/bal}
on CIFAR-10-LT against \R{t2/cifar10-r100/FedAvg/bal} for plain FedAvg and \R{t2/cifar10-r100/FedYoYo/bal} for
FedYoYo; its heads alone reach \R{t2/cifar10-r100/FedECS-LT (heads only)/bal}. The readout changes the heads by
\R{t2/cifar10-r100/FedECS-LT/gain_er} points on CIFAR-10-LT and \R{t2/cifar100-r100/FedECS-LT/gain_er} on
CIFAR-100-LT. The CIFAR-100-LT values of FedECS-LT are from one seed, because runs of FedECS-LT failed there where the
baselines finished; the missing runs are not random.

\paragraph{Why the readout does not help here} The readout repairs a last layer fitted to one client's prior
(\cref{sec:s-lastlayer}). With one shared model, the head is trained by every client in turn and averaged by FedAvg, so
its last layer already sees every client's classes. There is nothing left for the readout to repair, and a linear
discriminant on non-Gaussian embeddings can be worse than a trained last layer. The paper states that it did not test
why the readout lowers the accuracy.

\paragraph{Why the generator does not help here} With one shared model, FedAvg itself carries every client's
knowledge into every client's model, which is the job the generator does when models differ. FedECS-LT also trains its
head on frozen features of unaugmented images, while the baselines train the whole network end to end on augmented
images (\cref{ch:design}). The paper offers the second difference as a possible explanation and states that it did not
test it.

\paragraph{The published numbers} The baselines run in their authors' library on the paper's splits and report the
model after the last round. Against FedYoYo's own evaluation in the same protocol, most of them are within about three
points of the published values. Three (CReFF, Fed-GraB, FedLC) are further below; they run with the library's common
settings, which differ from those of their own papers, and were not tuned.

\paragraph{What it shows} FedECS-LT is a method for clients with different models. When the clients can share one
model, a plain FedAvg is better on CIFAR-10-LT, and the long-tail methods built for that setting are much better.

\subsection{Representations}

\paragraph{The question} Do the bodies of FedECS-LT separate the classes better than those of GeFL-F, independently of
any last layer?

\paragraph{The measurement} The leave-one-out 10-nearest-neighbour accuracy of the test embeddings classifies each
test image by the majority class of its ten nearest other test embeddings. It uses no last layer at all. It is
\R{repr/cifar10/FedECS-LT/knn} for FedECS-LT, \R{repr/cifar10/GeFL-F/knn} for GeFL-F and \R{repr/cifar10/FedDF/knn}
for FedDF on CIFAR-10-LT, averaged over the ten architectures and three seeds. FedECS-LT is higher in
\R{repr/cifar10/knn_wins/GeFL-F} of the \R{repr/cifar10/knn_cells} architecture--seed pairs against GeFL-F and in
\R{repr/cifar10/knn_wins/FedDF} against FedDF. \Cref{fig:s-tsne} shows a t-SNE map of the embeddings of one
architecture.

\begin{figure}[t]
\centering
\slotfig{fig_tsne}{wide}{t-SNE of GeFL-F and FedECS-LT embeddings.}
\caption{t-SNE of the test embeddings of one architecture on the same 2,000 CIFAR-10-LT test images, (a) GeFL-F and (b)
FedECS-LT; tail classes drawn larger. Look only at which points group together: in a t-SNE map the sizes of clusters
and the distances between them carry no meaning.}
\label{fig:s-tsne}
\end{figure}

\paragraph{What it shows and does not} The bodies of FedECS-LT separate the classes a little better, consistently
across architectures and seeds. The difference is a few points, which agrees with the small same-readout margin of
\cref{sec:s-same-readout}. A t-SNE map is an illustration; it cannot carry a quantitative claim.

\subsection{Calibration, the weakest architecture and cost}

\paragraph{Calibration} The expected calibration error of FedECS-LT is \R{main/cifar10/FedECS-LT/ece} on CIFAR-10-LT,
against \R{main/cifar10/GeFL-F/ece} for GeFL-F and \R{main/cifar10/strongest/ece} for \R{main/cifar10/strongest}; it
ranks \R{main/cifar10/ece_rank} of \R{main/cifar10/ece_n} methods there. On CIFAR-100-LT it ranks
\R{main/cifar100/ece_rank}, with \R{main/cifar100/FedECS-LT/ece} against \R{main/cifar100/GeFL-F/ece} for GeFL-F and
\R{main/cifar100/FedDF/ece} for FedDF. So calibration is not an advantage on CIFAR-100-LT.

\paragraph{The weakest architecture} Balanced accuracy averages ten models. The weakest architecture of FedECS-LT
reaches \R{main/cifar10/FedECS-LT/worst_arch} on CIFAR-10-LT and \R{main/cifar100/FedECS-LT/worst_arch} on
CIFAR-100-LT, so the mean does not hide a failed model.

\begin{table}[t]
\centering\small
\caption{Cost on CIFAR-10-LT, from the paper's cost table: communication per client over the whole training, wall-clock time
of a run (on a shared T4 GPU, so indicative only), forward MFLOPs per image and training GFLOPs per round. Compare the
first row with Balanced Softmax.}
\label{tab:s-cost}
\slottab{cost}{}
\end{table}

\paragraph{Cost} \Cref{tab:s-cost} gives the measured cost. FedECS-LT takes \R{cost/cifar10/time_ratio} times the
wall-clock time of GeFL-F and \R{cost/cifar10/comm_ratio} times its communication. Most of the extra time is the longer
synthetic phase ($T_s=10$), and most of the extra communication is the generator, which is larger than GeFL-F's
because of its projection of the class means. Against Balanced Softmax, which trains its heads locally, FedECS-LT
communicates \R{cost/cifar10/comm_ratio_bsm} times as much and takes \R{cost/cifar10/time_ratio_bsm} times as long,
for \R{main/cifar10/er_diff/Balanced Softmax} points when both have the readout. This is the trade a reader must weigh,
and \cref{ch:defence} returns to it.

\subsection{The experiments still pending}\label{sec:s-pending}

Five experiments asked for by the reviews were built as Flower notebooks and had not returned when this text was
written. The paper's text for each is already written and appears by itself once the results are filled in. For each,
here is the question, what result would support or weaken the paper, and what changes.

\paragraph{Several clients per architecture} The paper's 50-client setting has five clients per architecture, so heads
of one architecture are averaged in Stage III. The question: is the readout's dominance an artefact of one client per
architecture? With averaging, Balanced Softmax's heads would no longer be purely local, and their last layers would no
longer fit one client's prior, so the readout should matter less. The paper is supported if FedECS-LT stays ahead of
every method with the readout by more than the noise floor; it is weakened if the margins fall to zero, because then
FedECS-LT's case would rest on the one-client-per-architecture setting alone.
\numelse{k50/cifar10/margin}{\numelse{k50/cifar100/er_margin}{\emph{Result:} with 50 clients FedECS-LT is \R{k50/cifar10/margin} and
\R{k50/cifar100/margin} points ahead of the strongest other method without the readout, and
\R{k50/cifar10/er_margin} and \R{k50/cifar100/er_margin} with it.}{The runs are pending; the first attempt on
CIFAR-10-LT ran out of host memory on Kaggle (run log) and was rebuilt to run one job per GPU.}}{The runs are pending; the first attempt on
CIFAR-10-LT ran out of host memory on Kaggle (run log) and was rebuilt to run one job per GPU.}

\paragraph{New client splits} Seeds 6--8 draw splits that the development never used, for FedECS-LT, GeFL-F, FedDF
and Balanced Softmax. The question: are the CIFAR margins the product of choosing settings on the reported splits? The
paper is supported if the margins on seeds 6--8 are like those on 3--5, and weakened if the same-readout margin
vanishes there.
\numelse{fresh/cifar10/margin}{\numelse{fresh/cifar100/er_margin}{\emph{Result:} on the new splits FedECS-LT is \R{fresh/cifar10/margin} and
\R{fresh/cifar100/margin} points ahead without the readout, and \R{fresh/cifar10/er_margin} and
\R{fresh/cifar100/er_margin} with it.}{The runs are pending.}}{The runs are pending.}

\paragraph{Variants of GeFL-F} Seven variants isolate single changes, among them GeFL-F with decoupled weight decay
alone (AdamW) and GeFL-F and FedDF with ten synthetic epochs. The AdamW arm is the direct test of the decay diagnosis
of \cref{sec:s-decay}: if it raises GeFL-F's balanced accuracy or tail recall by more than the noise floor, the class
row shrinkage is shown to cost accuracy; if not, the shrinkage is a property of the generator without a measured
consequence. The ten-epoch arms test whether the heads' advantage is computation: if GeFL-F or FedDF with ten
synthetic epochs reach FedECS-LT's heads, it is.
\numelse{abl/cifar10/GeFL-F + AdamW/gain}{\numelse{abl/cifar100/GeFL-F + AdamW/gain}{\emph{Result:} decoupled decay alone changes GeFL-F by
\R{abl/cifar10/GeFL-F + AdamW/gain} points on CIFAR-10-LT and \R{abl/cifar100/GeFL-F + AdamW/gain} on
CIFAR-100-LT.}{The runs are pending.}}{The runs are pending.}

\paragraph{Components of FedECS-LT} Each run removes one component: the mean anchor, CCWA, the moment correction,
consensus soft labels, logit adjustment, or nine of the ten synthetic epochs. One more replaces the trained generator
by Gaussians with the exact class means and covariances, which tests whether a trained generator is needed at all. The
question is RQ5. Each component is shown to matter if removing it lowers balanced accuracy with the readout by more
than \R{repeat/er/floor} points, significantly. If all removals stay within the floor, no component has a measured
effect beyond the readout, and the paper must say so. If the Gaussian sampler matches the full method, Stage II is
unnecessary.
\numelse{comp/cifar10/largest}{\emph{Result:} with the readout, the largest loss on CIFAR-10-LT comes from
``\R{comp/cifar10/largest}'' (\R{comp/cifar10/largest_drop} points), and \R{comp/cifar10/n_small} of the
\R{comp/n_variants} changes stay within the noise floor.}{The runs are pending.}

\paragraph{Other settings} Sweeps over the label skew ($\alpha=0.1$ and $1.0$), the number of clients (50 and 100)
and the imbalance ratio (50 and 10) on CIFAR-10-LT test whether the result depends on the main setting's values. Their
runs are pending, and the paper's text for them is not yet written.

What, taken together, does this evidence support?
